\input{../tex-inputs/latex-leseansicht-vorspann}

\section[Arthur und Olga Schnitzler an Gustav Schwarzkopf, 20. 8. 1914]{L04488 Arthur und Olga Schnitzler an Gustav Schwarzkopf, 20. 8. 1914}
\nopagebreak\mylabel{L04488v}
\rehead{ }\normalsize\beginnumbering\briefempfaengerindex{Schwarzkopf, Gustav@\textsc{Schwarzkopf, Gustav}!zzzSchnitzler, Olga@\emph{von Olga Schnitzler}!1914-08-202@{20. 8. 1914}|(be}\briefempfaengerindex{Schwarzkopf, Gustav@\textsc{Schwarzkopf, Gustav}!zzzSchnitzler, Arthur@\emph{von Arthur Schnitzler}!1914-08-202@{20. 8. 1914}|(be}
\toendnotes[C]{\smallbreak\pagebreak[2]}
\correspDesc{Versand  durch Arthur Schnitzler, Olga Schnitzler am 20. 8. 1914 in Salzburg
\newline{}Übermittlung  am 20. 8. 1914 in Salzburg
\newline{}Erhalt  durch Gustav Schwarzkopf im Zeitraum [21. 8. 1914 – 25. 8. 1914?] in Wien}\toendnotes[C]{\smallbreak}
\Standort{DLA, A:Schnitzler, HS.1985.1.1897.}
\physDesc{Bildpostkarte, 345 Zeichen
\newline{}Handschrift Arthur Schnitzler: Bleistift, deutsche Kurrent
\newline{}Handschrift Olga Schnitzler: Bleistift, lateinische Kurrent
\newline{}Versand: Stempel: »\nobreak{}\oindex{Salzburg@\textbf{Salzburg}|pwk}Salzburg 2, 20. VIII. 14, 9\nobreak{}«.  }
\pstart
           \noindent{}{\pb}\textcolor{gray}{\textbf{Salzburg. Staatsbrücke\oindex{Staatsbrücke@\textbf{Staatsbrücke}|pw}.}}\pend
           \pstart{}{\pb}Hrn. \textsc{Gustav Schwarzkopf}\pend{}\pstart{}Wien I\oindex{I., Innere Stadt@\textbf{I., Innere Stadt}|pw}\pend{}\pstart{}\textsc{Tiefer Graben 17}\oindex{Wien@\textbf{Wien}!I., Innere Stadt@\textbf{I., Innere Stadt}!Tiefer Graben 17@\textbf{Tiefer Graben 17}|pw}\pend{}{\bigskip}\vspace{1em}
\pstart
           \noindent{}Wieder einmal hier und bei
               herrlichſtem Wetter. Wir{ }ſind
               von Buchs\oindex{Buchs@\textbf{Buchs}|pw} bis I{\geminationn}sbruck\oindex{Innsbruck@\textbf{Innsbruck}|pw} 14 Stunden,
               von I{\geminationn}sbruck\oindex{Innsbruck@\textbf{Innsbruck}|pw} nach Salzburg\oindex{Salzburg@\textbf{Salzburg}|pw} 21
               gefahren! – Morgen nach Iſchl\oindex{Bad Ischl@\textbf{Bad Ischl}|pw}, für
               ein paar Tage; – dann, hoffentlich 
               in 2 Tagesreiſen nach Wien\oindex{Wien@\textbf{Wien}|pw}.\pend
           
\pstart
           Auf Wiederſehen – alles ubrg
               auf mündlich! Herzlichſt Ihr{\\[\baselineskip]}\spacefill\mbox{Arthur}\pend
           \leftskip=0em{}\selectlanguage{ngerman}\vspace{1em}
\pstart
           \noindent{}{\pb}{[}hs. O. Schnitzler:{]} Alles Herzliche von Ihrer \spacefill\mbox{O. S.}\pend
           \selectlanguage{ngerman}\endnumbering\briefempfaengerindex{Schwarzkopf, Gustav@\textsc{Schwarzkopf, Gustav}!zzzSchnitzler, Olga@\emph{von Olga Schnitzler}!1914-08-202@{20. 8. 1914}|)be}\briefempfaengerindex{Schwarzkopf, Gustav@\textsc{Schwarzkopf, Gustav}!zzzSchnitzler, Arthur@\emph{von Arthur Schnitzler}!1914-08-202@{20. 8. 1914}|)be}\mylabel{L04488h}  \newcommand{\dateiname}{L04488}\newcommand{\titel}{Arthur und Olga Schnitzler an Gustav Schwarzkopf, 20. 8. 1914}\newcommand{\editorInnen}{Herausgegeben von Jahnke, SelmaMüller, Martin Anton}\input{../tex-inputs/latex-leseansicht-abspann}
